\begingroup
\setlength{\parskip}{2pt}
\tcbset{before skip=4pt,after skip=4pt,top=4pt,bottom=4pt}
\source{Лекции, с. 41–45; схема через цикл и его части; минимальность и локальные конфигурации раскрыты составителем}
\begin{definition}[Топологический подграф]
Подграф, гомеоморфный \(H\), — подразбиение \(H\): его рёбра заменены
простыми путями с различными новыми внутренними вершинами.
Вершины, соответствующие вершинам \(H\), называются главными.
\end{definition}
\begin{theorem}[Понтрягин–Куратовский]
Граф планарен тогда и только тогда, когда не содержит подграфа,
гомеоморфного \(K_5\) или \(K_{3,3}\).
\end{theorem}
\begin{proof}[Необходимость]
Подграф планарного графа планарен, подразделение и обратное склеивание
рёбер сохраняют планарность, а \(K_5\), \(K_{3,3}\) непланарны по билету 15.
\end{proof}
\subsection*{Лекционная схема достаточности}
В вопросе требуется \emph{схема}. Ниже сохранён ход лекций:
минимальный контрпример, цикл, внутренние и внешние части, запрещённые
конфигурации. Геометрическая редукция леммы 14 (с. 44–45) остаётся
схематической; определения, предварительные шаги и извлечение двух
стандартных конфигураций расписаны явно. Теорема о стягиваемом ребре
и другое книжное доказательство здесь не используются.

\textbf{1. Минимальный контрпример и трёхсвязность.}
Выберем непланарный граф \(G\) без двух запрещённых подграфов,
минимальный сначала по \(n\), затем по \(m\).
Любой его собственный подграф планарен.
Граф связен и не имеет точек сочленения: иначе меньшие части планарны
и их укладки склеиваются по билету 16. При \(n\ge3\) мост также
давал бы точку сочленения. Малые графы планарны, поэтому они не контрпримеры.

Покажем, что нет и разделяющей пары \(\{a,b\}\).
Разобьём компоненты \(G-\{a,b\}\) на две непустые группы и добавим
к каждой \(a,b\), получив \(G_1,G_2\). Каждый компонент смежен и с
\(a\), и с \(b\), иначе одна из этих вершин была бы точкой сочленения.
Значит, в каждой части есть \(ab\)-путь с внутренними вершинами в этой части.
Рассмотрим \(G_i+ab\). Если в нём есть запрещённое подразбиение,
оно использует новое ребро \(ab\); заменим его путём через другую часть.
Получим запрещённое подразбиение в \(G\), противоречие.
По минимальности обе дополненные части планарны. Выносим \(ab\)
на внешнюю грань и склеиваем укладки по ребру с разных сторон;
если ребра не было, удаляем его. Получаем планарность \(G\), противоречие.
Таким образом, \(G\) трёхсвязен: удаление любых одной или двух вершин
сохраняет связность.

\textbf{2. Удалённое ребро и цикл (лемма 13 лекций).}
Выберем \(uv\in E(G)\), положим \(H=G-uv\). Граф \(H\) планарен
по минимальности и двусвязен. Действительно, для любой вершины \(z\)
граф \(G-z\) двусвязен по трёхсвязности \(G\), поэтому ребро \(uv\),
если оно осталось, не мост в \(G-z\). Значит, \(H-z\) связен.
По билету 5 в \(H\) есть простой цикл \(C\) через \(u,v\).

\begin{definition}[Части относительно цикла]
Частью относительно \(C\) называется либо хорда \(C\), либо компонента
\(H-V(C)\) вместе со всеми рёбрами к \(C\) и их концами на \(C\).
Эти концы — точки прикрепления. В плоской укладке внутренности каждой
части лежат по одну сторону \(C\). Обозначаем внешние части \(A_j\),
внутренние — \(B_i\), как в лекциях. Это \emph{не} мосты-рёбра из билета 4.
\end{definition}
На сфере можно выбирать укладку \(H\), цикл \(C\) через \(u,v\) и его
сторону так, чтобы на выбранной внутренней стороне было как можно больше
граней. Максимум существует: число граней конечно и фиксировано формулой Эйлера.
Именно такой экстремальный выбор позволяет перестраивать цикл,
сравнивая число захваченных граней (рисунок с. 42–43).

Если внешняя часть имеет хотя бы три точки прикрепления, внутри неё
можно выбрать дерево с тремя выходами на \(C\). Между тремя выходами
есть дуга \(C\), внутренность которой не содержит ни \(u\), ни \(v\).
Заменяя её путём по двум ветвям дерева, выбираем новую внутреннюю сторону,
содержащую прежнюю и ещё область между дугой и деревом. Одна из двух
остальных областей остаётся снаружи. Число внутренних граней растёт,
противоречие максимальности. В связном графе без точек сочленения у
части не меньше двух точек прикрепления. Поэтому у каждой \(A_j\)
ровно две. Внешняя часть не может иметь внутренних вершин: её две
точки прикрепления разделяли бы \(G\). Следовательно, все \(A_j\) — хорды.

\begin{center}
\begin{tikzpicture}[scale=.65,every node/.style={circle,fill=cbBlue,inner sep=2pt}]
\draw (0,0) circle (1.5);
\node (u) at (0,1.5) {};\node[fill=none,above] at (u) {$u$};
\node (v) at (0,-1.5) {};\node[fill=none,below] at (v) {$v$};
\node (a) at (1.5,0) {};\node[fill=none,right] at (a) {$a$};
\node (b) at (-1.5,0) {};\node[fill=none,left] at (b) {$b$};
\draw[cbOrange,thick] (a) .. controls (3,3.3) and (-3,3.3) .. (b);
\node[fill=none] at (1.6,2.1) {$A_j$};
\node (z) at (0,0) {};
\draw[cbGreen,thick] (z)--(u);\draw[cbGreen,thick] (z)--(v);
\draw[cbGreen,thick] (z)--(a);\draw[cbGreen,thick] (z)--(b);
\node[fill=none] at (.65,.5) {$B_i$};\node[fill=none] at (-1.1,1.2) {$C$};
\end{tikzpicture}
\end{center}
Внешняя хорда \(A_j\) и внутренняя часть \(B_i\) прикреплены к одному
циклу \(C\). Само ребро \(uv\) удалено; задача — вернуть его без пересечений.

\textbf{3. Разделение частей (лемма 14 лекций).}
Две пары точек \(a,b\) и \(c,d\) разделяют друг друга на цикле,
если четыре точки различны и встречаются в порядке \(a,c,b,d\).
Пути между такими парами не могут лежать по одну сторону \(C\):
путь первой пары вместе с дугой цикла отделяет концы второго по теореме Жордана.
Более общо, две части конфликтуют, если имеют чередующиеся пары точек
прикрепления или не менее трёх общих точек прикрепления.
Для хорды \(ab\) конфликт с частью \(B\) означает наличие у \(B\)
точек в обеих открытых дугах \(C-\{a,b\}\).

Ребро \(uv\) тоже нужно разместить по одну из сторон цикла.
В лекционной геометрической редукции сравниваются перестановки частей
между сторонами и замены дуг \(C\) путями через части. Если \(uv\)
разместить нельзя, выделяется конфликт внутренних и внешних частей,
разделяющий также \(u,v\); сокращение несущественных ветвей даёт
одну из запрещённых конфигураций ниже. Если конфликта нет, части
переносятся на подходящие стороны и дуга \(uv\) дорисовывается —
получается укладка \(G\), противоречие.

\textit{Статус этого шага.} Полный перебор взаимного расположения точек
прикрепления здесь, как и в лекционной схеме, не доказывается.
В частности, произвольный конфликт двух частей сам по себе ещё не
доказывает непланарность: требуется невозможность распределить все части
и \(uv\) по сторонам. Рисунки ниже объясняют заключительные конфигурации,
а не заменяют эту редукцию необоснованным утверждением о любом конфликте.

\textbf{4. Как возникают запрещённые подграфы.}
\emph{Конфигурация \(K_{3,3}\).} На цикле по порядку лежат
\(u,a,c,v,b,d\), а вне его выбраны три пути \(uv,ab,cd\)
с попарно непересекающимися внутренностями, не встречающими \(C\).
Сохраним шесть дуг между соседними отмеченными вершинами и эти три пути.
Это подразбиение \(K_{3,3}\) с долями \(\{u,c,b\}\), \(\{a,v,d\}\).
Каждая главная вершина соединена с тремя вершинами другой доли;
все девять путей имеют разные внутренние вершины.

\emph{Конфигурация \(K_5\).} На цикле по порядку лежат \(u,a,v,b\),
имеются противоположные пути \(uv,ab\), а внутренняя часть содержит
вершину \(z\) с четырьмя внутренне непересекающимися выходами к
\(u,a,v,b\). Сохраняем четыре дуги цикла, два противоположных пути
и четыре выхода: получаем десять путей подразбиения \(K_5\)
с главными вершинами \(u,a,v,b,z\).
В геометрическом разборе выбираются именно непересекающиеся ветви;
при нескольких развилках выделяются соответствующие пути и удаляется лишнее.

\begin{center}
\begin{tikzpicture}[scale=.8,every node/.style={circle,fill=cbBlue,inner sep=2pt}]
\foreach \i/\lab in {1/u,2/a,3/c,4/v,5/b,6/d}{
 \node (p\i) at ({90-60*(\i-1)}:1.25) {};
 \node[fill=none] at ({90-60*(\i-1)}:1.55) {$\lab$};}
\draw (p1)--(p2)--(p3)--(p4)--(p5)--(p6)--(p1);
\draw[cbOrange,thick] (p1)--(p4);\draw[cbGreen,thick] (p2)--(p5);
\draw[cbPurple,thick] (p3)--(p6);
\node[rectangle,fill=none] at (0,-2) {$K_{3,3}$};
\begin{scope}[xshift=5cm]
\foreach \i/\lab in {1/u,2/a,3/v,4/b}{
 \node (q\i) at ({90-90*(\i-1)}:1.25) {};
 \node[fill=none] at ({90-90*(\i-1)}:1.55) {$\lab$};}
\node (z) at (0,0) {};\node[fill=none,above right] at (z) {$z$};
\draw (q1)--(q2)--(q3)--(q4)--(q1);
\foreach \i in {1,...,4}{\draw (z)--(q\i);}
\draw[cbOrange,thick] (q1) to[bend left=28] (q3);
\draw[cbGreen,thick] (q2) to[bend left=35] (q4);
\node[rectangle,fill=none] at (0,-2) {$K_5$};
\end{scope}
\end{tikzpicture}
\end{center}
Это схемы подграфов, а не плоские укладки: пересечения линий без отмеченной
вершины не означают общей вершины соответствующих путей.

Итак, геометрический шаг либо даёт укладку \(G\), либо выделяет
запрещённое подразбиение. Оба исхода противоречат выбору контрпримера.
Это завершает \emph{схему} достаточности в форме лекций.

\endgroup
